\section{Weights for flat bundles of algebraic families}
\label{sec:diagram}

Consider first a product $S\times X^{\mathrm B}$, with $S$ a closed
oriented smooth surface.  The moving-point construction in the proof overview
requires integration over $S$ to lower cohomological degree by two while
preserving weights.  On this product, give the K\"unneth factor
$H^a(S,\mathbb Q)$ pure weight zero and Hodge type $(0,0)$ in every degree,
and retain Deligne's mixed Hodge structure on $H^*(X^{\mathrm B},\mathbb Q)$.
Integration over $S$ then acts on the first factor without a Tate twist.
Even when $S$ is an algebraic curve, this convention differs from its
ordinary Deligne weights.

The moving-puncture flag spaces need not be products.  We therefore
extend this convention to bundles whose transition maps are locally
constant algebraic automorphisms of a fiber family.  A finite diagram of
algebraic fibers will give the mixed Hodge structures together with their
compatibility with parallel transport and proper Gysin maps.  Section~\ref{sec:logarithmic}
will construct the required bundles and identify the scalar unflagged
bundle as an actual flat product.  The simplicial mixed-Hodge formalism
goes back to Deligne \cite[Section~8.3]{Deligne1974}; Shende uses weight-zero
simplicial cochains to compute weights of tautological classes on
character varieties \cite[proof of the main Theorem]{Shende2017}.
Here we supply the relative construction and the Gysin compatibility
needed for the moving-point correspondence.

\subsection{The finite diagram}

Let $V$ be a smooth connected complex algebraic variety and let
$p:H\to V$ be an algebraic morphism of finite type.  We assume that
$Rp_*\mathbb Z_H$ has bounded, locally constant, finitely generated
cohomology sheaves and that its natural ordinary base-change maps
\begin{equation}
 \label{diag:basechange}
 i_t^*Rp_*\mathbb Z_H\longrightarrow
 R\Gamma(H_t,\mathbb Z)
\end{equation}
are isomorphisms for all $t\in V$, where $i_t:\{t\}\hookrightarrow V$.
These are hypotheses about ordinary direct image; smoothness alone would
not imply them for a nonproper morphism.

Let $S$ be a compact smooth manifold.  A \emph{flat bundle of the family
$H\to V$ over $S$} means a space $\mathscr H$ with a map to $S\times V$
and local identifications
\[
 \mathscr H|_{U_i\times V}\simeq U_i\times H
\]
whose transition maps have the form
$(s,h)\mapsto(s,g_{ij}(h))$, with $g_{ij}$ an algebraic automorphism
over $V$, locally constant in $s$.  The maps $g_{ij}$ satisfy the cocycle
identity.  Write $\mathscr H_t$ for the total space over $S$ at $t\in V$.
Thus $\mathscr H_t$ includes the parameter space $S$; it need not be an
algebraic variety.

A map of flat bundles will mean one which, in simultaneous flat
trivializations, is the identity on $S$ times an algebraic map of the
fiber families.

Choose a finite good cover $\mathcal U=\{U_i\}$ refining a flat
trivializing cover.  Every nonempty intersection $U_I=\bigcap_{i\in I}U_i$
is contractible.  Choose a trivialization on each $U_I$.  Inclusion of
intersections then gives a finite diagram of copies
$p_I:H_I\to V$ and algebraic isomorphisms over $V$.
The cocycle identity makes this a diagram of maps of varieties, before
passing to cohomology.  Its index category is the poset of nonempty
intersections, ordered in the direction of cohomological restriction.

We use the enhanced bounded derived category of mixed Hodge modules,
whose six operations and rational realization are compatible with
diagrams \cite[Section~2.1, Theorem~2.1.1, Remark~2.1.2, and
Proposition~4.2.6]{TubachStacks2025}.
For an algebraic variety $T$, let $\mathbb Q_T^H$ denote the constant
mixed-Hodge-module complex with underlying complex $\mathbb Q_T$,
without a dimension shift.  Define
\begin{equation}
 \label{diag:complex}
 K_{\mathscr H}
   =\mathop{\mathrm{holim}}_I Rp_{I*}\mathbb Q_{H_I}^H
   \quad\text{in }D^b\mathrm{MHM}(V).
\end{equation}
Here the homotopy limit means the derived limit, retaining the
transition maps and their compatibilities.  The index poset has a
finite-dimensional nerve, so this is a finite limit and remains
bounded.  In particular, \eqref{diag:complex} uses the full derived
diagram, not only the cohomology groups of its terms.

\begin{proposition}[Diagram weights]
 \label{diag:variation}
 Under the preceding hypotheses, $H^j(\mathscr H_t,\mathbb Q)$ carries
 a graded-polarizable mixed Hodge structure, functorial for maps of flat
 bundles that are algebraic in flat trivializations.  As $t$ varies,
 these structures form an admissible graded-polarizable variation of
 mixed Hodge structures, with an integral underlying local system up
 to torsion.  They are computed by $H^j(i_t^*K_{\mathscr H})$.

 If $H^j(H_t,\mathbb Q)$ is Hodge--Tate for every $j$ and $t$, then
 $H^j(\mathscr H_t,\mathbb Q)$ is Hodge--Tate as well.  We call its
 weight filtration the \emph{diagram weight filtration}.
\end{proposition}

\begin{proof}
 First we identify the underlying ordinary complex.  Open-cover
 descent computes the ordinary direct image of $\mathscr H\to V$
 from the spaces $U_I\times H_I$.  Projection to $H_I$ induces a
 quasi-isomorphism on direct images over $V$: over any open subset
 of $V$, its fiber in this product is the contractible space $U_I$.
 These projections commute with the transition maps.  Their derived
 limit therefore identifies the rational realization of
 $K_{\mathscr H}$ with the ordinary direct image of $\mathscr H\to V$.
 The same argument works with integral coefficients.

 A finite limit of bounded complexes with locally constant cohomology
 again has locally constant cohomology.  Its integral cohomology
 sheaves have finite type.  Point pullback commutes with finite limits,
 and the termwise base-change maps are isomorphisms by
 \eqref{diag:basechange}.  The corresponding maps in mixed Hodge
 modules are isomorphisms because rational realization is conservative.
 It follows that $i_t^*K_{\mathscr H}$ computes the derived diagram of
 the algebraic fibers $H_{I,t}$ and that its underlying cohomology is
 $H^*(\mathscr H_t,\mathbb Q)$.

 Put $v=\dim V$.  Since the ordinary cohomology sheaves of
 $K_{\mathscr H}$ are locally constant, its perverse cohomology
 ${}^pH^{j+v}(K_{\mathscr H})$ is a smooth mixed Hodge module.
 Smooth mixed Hodge modules, after the dimension shift, are admissible
 graded-polarizable variations
 \cite[Section~3.2, Equation~(3.2.1)]{SaitoExtension1990}.
 Applying point pullback to the perverse truncations shows that the
 fiber of this variation is $H^j(i_t^*K_{\mathscr H})$: a smooth
 perverse object has only one nonzero fiber degree, namely $-v$.

 Finally, mixed Hodge structures whose weight gradeds have only Hodge
 types $(r,r)$ are closed under subobjects, quotients, and extensions.
 The long exact sequence of a cone thus shows that complexes with
 Hodge--Tate cohomology are closed under finite limits.  Applying this
 observation to the fiber diagram proves the last assertion.
 Algebraic maps of the diagrams give the stated functoriality.
 A common good refinement compares different choices of cover and
 trivializations: its comparison maps are isomorphisms after ordinary
 descent, hence also in mixed Hodge modules by conservativity.
\end{proof}

The construction also preserves the cup product.  Each algebraic
direct image of the constant complex has its functorial cup product,
and the transition pullbacks preserve it.  The derived limit inherits
this product.  Under the open-cover comparison in the proof of
Proposition~\ref{diag:variation}, it is the ordinary cup product on
$H^*(\mathscr H_t,\mathbb Q)$.  Thus parallel transport in $V$ preserves
both the cohomology ring and its weight filtration.

\subsection{Products and proper Gysin maps}

We next identify the maps that preserve these weights.

\begin{proposition}[Weight shifts]
 \label{diag:gysin}
 The diagram weight filtrations have the following properties.
 \begin{enumerate}
 \item For the product $S\times M$, with $M$ a complex algebraic
 variety, the ordinary K\"unneth isomorphism identifies its mixed Hodge
 structure with that on
 \[
 H^*(S,\mathbb Q)\otimes H^*(M,\mathbb Q),
 \]
 where every $H^a(S,\mathbb Q)$ is pure of weight zero and type $(0,0)$,
 and $H^*(M,\mathbb Q)$ has its Deligne mixed Hodge structure.
 If $S$ is oriented and closed of real dimension $d_S$, projection
 pullback and integration over $S$ preserve weight; the latter has
 cohomological degree $-d_S$.
 \item Pullbacks by maps of flat bundles preserve weight.
 \item Suppose $u:\mathscr H_t\to\mathscr H'_t$ is a proper map of
 flat bundles with smooth equidimensional algebraic fibers.  Put
 $c=\dim_{\mathbb C}H'_t-\dim_{\mathbb C}H_t$.  With the complex
 orientations on the fibers, its ordinary Gysin map satisfies
 \[
 u_*:W_aH^j(\mathscr H_t)\longrightarrow
       W_{a+2c}H^{j+2c}(\mathscr H'_t).
 \]
 If $S$ is oriented, this is the Gysin map for the induced orientations
 on the total spaces.
 \end{enumerate}
\end{proposition}

\begin{proof}
 For a product, the diagram is constant.  Its derived limit is
 \[
 C^*(N\mathcal U,\mathbb Q(0))\otimes R\Gamma(M,\mathbb Q_M^H),
 \]
 where $N\mathcal U$ is the finite nerve of the good cover.  The
 first complex has a copy of $\mathbb Q(0)$ for each simplex, with
 its ordinary simplicial differential.  All its cohomology therefore
 has weight zero.  The external-product comparison on the cover
 gives exactly the ordinary K\"unneth map.  Evaluation on a simplicial
 fundamental cycle defines
 \[
 C^*(N\mathcal U,\mathbb Q(0))\longrightarrow
       \mathbb Q(0)[-d_S].
 \]
 This is a morphism of complexes of mixed Hodge structures and
 realizes ordinary integration over $S$.  It has no Tate twist.
 Projection pullback is induced by the unit of the same cochain
 complex.  This proves the first assertion; functoriality of algebraic
 pullbacks proves the second.

 For the third assertion, fix $t$ and choose simultaneous
 trivializations of the map.  On each algebraic fiber map
 $u_{I,t}:H_{I,t}\to H'_{I,t}$, smooth
 dualizing complexes and the proper counit give the Gysin morphism
 \begin{equation}
  \label{diag:trace}
  Ru_{I,t*}\mathbb Q^H_{H_{I,t}}\longrightarrow
       \mathbb Q^H_{H'_{I,t}}[2c](c).
 \end{equation}
 Its degree and Tate twist give the displayed weight shift.
 Complex orientations are preserved by the algebraic transition
 isomorphisms, and the proper counit is natural.  Consequently
 \eqref{diag:trace} gives a coherent map of the finite diagrams and
 hence of their derived limits.

 It remains to identify its ordinary realization.  On
 $U_I\times H_{I,t}$ the map is the identity on $U_I$ times $u_{I,t}$.
 The orientation trace
 for this product is the fiber trace, tensored with the identity on
 $U_I$.  It is compatible with restriction to intersections and with
 the transition isomorphisms.  Open-cover descent therefore identifies
 the realized map with the ordinary topological Gysin map of the
 total spaces.  In particular, the construction gives the actual
 Gysin map, not merely a map with the same source and target.
\end{proof}

\subsection{Finite monodromy on weight gradeds}

The diagram construction has now supplied a variation on the
cohomology of the total spaces.  Hodge--Tate fibers and the integral
structure put a strong restriction on its monodromy.

\begin{lemma}
 \label{diag:finite}
 Assume the Hodge--Tate hypothesis of
 Proposition~\ref{diag:variation}.  The combined monodromy action
 on all $\operatorname{Gr}^W_aH^j(\mathscr H_t,\mathbb Q)$ has finite
 image.  For any loop with cohomology monodromy $T$, there is an integer
 $M>0$ such that $T^M$ is unipotent and
 \begin{equation}
  \label{diag:log}
  N=\frac{1}{M}\log(T^M)
  \quad\text{satisfies}\quad N(W_a)\subset W_{a-2}.
 \end{equation}
 The logarithm is a finite polynomial.  On the total cohomology ring,
 $N$ is a derivation of cohomological degree zero.
\end{lemma}

\begin{proof}
 Each nonzero weight graded is a polarizable variation of pure type
 $(r,r)$ and weight $2r$.  After the conventional sign, its flat
 polarization is positive definite on the underlying real vector
 space.  Let $\mathcal L$ be the torsion-free quotient of the integral
 cohomology local system.  Intersecting $\mathcal L$ with the rational
 weight sub-local systems and taking successive quotients gives an
 invariant full lattice in each rational graded piece.  These
 intersections are saturated because the subspaces are rational.

 The automorphisms of a lattice preserving a positive definite form
 constitute a finite group: in a fixed lattice basis, every column of
 such an automorphism has a prescribed bounded norm and hence lies
 in a finite set.  Thus each graded monodromy image is finite.  Only
 finitely many cohomological degrees and weights occur, which proves
 the combined assertion.

 Choose $M$ killing the action of the specified loop on all weight
 gradeds.  Then
 \[
 (T^M-1)W_a\subset W_{a-1}=W_{a-2}
 \]
 whenever $a$ is even; the identical conclusion for odd $a$ follows
 because successive odd and even steps coincide.  The operator
 $T^M-1$ is nilpotent, since the weight filtration is finite.
 Its finite logarithm has the same lowering bound.

 Parallel transport preserves cup products.  For every integer
 $k\geq0$, the polynomial operator $\exp(kMN)=T^{Mk}$ is therefore
 a ring automorphism.  For cohomology classes $\alpha,\beta$, the
 identity
 \[
 \exp(kMN)(\alpha\beta)
   =\exp(kMN)(\alpha)\exp(kMN)(\beta)
 \]
 is an equality of polynomials in $k$.  Comparing linear coefficients
 proves $N(\alpha\beta)=N(\alpha)\beta+\alpha N(\beta)$.
 Monodromy preserves cohomological degree, so $N$ does as well.
\end{proof}

To apply these results, we will verify locally constant algebraic
transition maps for the family with a moving marked point, together
with simultaneous such transitions for its flag correspondence.
We will also identify the scalar unflagged bundle as a flat product.
Once those inputs are established, Proposition~\ref{diag:gysin}
controls its push-pull maps, and Lemma~\ref{diag:finite} supplies the
weight-lowering logarithm on the same ordinary cohomology groups.
